% Agent fragment. All finding IDs begin PL-. Bibliography locators are inline.
\section{The plethystic bridge and the actual content of the archive}
\label{sec:plethysm-bridge}

\subsection{Scope, conventions, and known dictionary}
The archive evidence here consists of the supplied transcriptions, including their
deletions and uncertainty marks, not images of every manuscript page. A claim made
by the supplied study is not thereby a checked claim of this paper. In particular,
the study's order-two check of cote 133, scan 19 fails equivariantly, as proved
below. No assertion of historical priority is made for the elementary consequences
labelled PL below.

We use ``pre-$\lambda$-ring'' for an additive-to-multiplicative map
$\lambda_t:R\longrightarrow 1+tR[[t]]$ with first coefficient $x$, and
``$\lambda$-ring'' for one also satisfying $\lambda_t(1)=1+t$ and the
universal product and iteration identities. This is the convention of SGA 6, Expos\'e V,
Definitions 2.1 and 2.4; some authors call the latter a special $\lambda$-ring.
The vector-bundle Grothendieck ring of a quasi-compact scheme has the latter
structure, with $\lambda^i[E]=[\bigwedge^iE]$ (SGA 6, Expos\'e VI,
Theorem 3.2). The free $\lambda$-ring on one generator is the symmetric-function
ring $\Lambda=\mathbb Z[e_1,e_2,\ldots]$; see the attached Borger,
\emph{$\Lambda$-rings and the field with one element}, \S1.1, the paragraph
beginning ``The functor $W^*$ has a right adjoint''. The universal map to a
$\lambda$-ring $R$ sending the generator to $x$ evaluates $e_i$ as $\lambda^i(x)$.

The characteristic-zero dictionary is exact:
\[
 R=\bigoplus_{n\geq0}K_0(\operatorname{Rep}_{\mathbb C}(\mathfrak S_n))
 \xrightarrow{\operatorname{ch}}\Lambda,
 \qquad [\lambda]\longmapsto s_\lambda,
 \quad [\mathbf1_{\mathfrak S_n}]\longmapsto h_n,
 \quad [\operatorname{sgn}_n]\longmapsto e_n.
\]
Here the multiplication on $R$ is induction from
$\mathfrak S_a\times\mathfrak S_b$ to $\mathfrak S_{a+b}$. It is not the
inner tensor product. For the latter, within one degree,
\[
 \operatorname{ch}([\mu]\otimes[\nu])=s_\mu*s_\nu
   =\sum_{\lambda\vdash n}g(\lambda,\mu,\nu)s_\lambda.
\]
For composition of polynomial Schur functors,
\[
 \operatorname{ch}(S_\alpha\circ S_\beta)=s_\alpha[s_\beta]
   =\sum_{\lambda\vdash |\alpha||\beta|}p^\lambda_{\alpha,\beta}s_\lambda.
\]
The first multiplication adds degrees, the inner tensor product preserves the
common degree, and plethysm multiplies degrees. Exact standard locators are
Macdonald, \emph{Symmetric Functions and Hall Polynomials}, second edition,
I~(7.3), I~(7.9), and Appendix A~(7.3)--(7.4); the last formula already proves
ordinary plethysm coefficients nonnegative. The positive categorical structure
is made explicit in Baez--Moeller--Trimble, \emph{Schur Functors and Categorified
Plethysm}, Theorem 6.11. None of these statements supplies a polynomial-time
algorithm or an Ehrhart model.

For a finite group $G$, let $B(G)$ be its Burnside ring. There are two different
operations to distinguish:
\[
 \epsilon^i(X)=[\{A\subset X:|A|=i\}],\qquad
 \sigma^i(X)=[\operatorname{Sym}^iX].
\]
Both extend to pre-$\lambda$ structures (with the appropriate generating series)
by disjoint-union decompositions. Define the exterior-compatible structure
$\lambda^{\mathrm{op}}$ instead by
\[
 \lambda^{\mathrm{op}}_t(x)=\sigma_{-t}(x)^{-1}.
\]
This generally differs from $\epsilon$. Linearization
$L:B(G)\to R_{\mathbb C}(G)$, $[X]\mapsto[\mathbb C[X]]$, respects
$\sigma$ and $\lambda^{\mathrm{op}}$, but not $\epsilon$ in general.
These conventions agree with Bouc--R\"okaeus,
\emph{A note on the $\lambda$-structure on the Burnside ring},
\S1 and Theorem 1.1, equation (1.2). The following proofs isolate the relevant
small counterexamples rather than reprove that general theorem.

\subsection{PL-01: a false identity in the supplied reading of scan 19}
\paragraph{Statement PL-01.}
The identity $2\epsilon^2(X)+X=X^2$ is false in $B(G)$.
For $G=C_2$ and its free two-point orbit $E$, put $x=[E]$. Then
\[
 x^2=2x,\qquad \epsilon^2(x)=1,\qquad
 \sigma^2(x)=x+1,\qquad \lambda^{\mathrm{op},2}(x)=x-1.
\]
Thus neither the subset operation nor the exterior-compatible operation makes
all three order-two equalities transcribed on scan 19 true at once. The
cardinality equality remains true.

\begin{proof}
The diagonal and off-diagonal pairs in $E\times E$ are two free orbits, proving
$x^2=2x$. Its unique two-element subset is fixed, proving $\epsilon^2(x)=1$.
The two repeated multisets form a free orbit and the multiset containing both
elements is fixed; hence $\sigma^2(x)=x+1$. The coefficient of $t^2$ in
$\sigma_t(x)\lambda^{\mathrm{op}}_{-t}(x)=1$ gives
$\lambda^{\mathrm{op},2}(x)=x^2-\sigma^2(x)=x-1$.
The fixed-point homomorphism $\Phi_G:B(G)\to\mathbb Z$ takes $x$ to zero
and $1$ to one. Therefore it takes $2\epsilon^2(x)+x$ to two and $x^2$ to
zero, so these classes differ. It also distinguishes
$2\lambda^{\mathrm{op},2}(x)+x=3x-2$ from $2x$.
Applying cardinality instead gives $2\binom{2}{2}+2=2^2$.
\end{proof}

The supplied \texttt{SA17\_STUDY\_cote133\_lambda.md}, section
``Scan 19: $\lambda$- and $\sigma$-operations'', claims that with subsets and
multisets the two descriptions agree. PL-01 refutes precisely that assertion.
The correction preserves the true identity
$\operatorname{Sym}^2X\simeq\epsilon^2X\sqcup X$, and the true opposite-series
identity, but uses different symbols for their two exterior-like operations.

\subsection{PL-02: what survives linearization, and why specialness matters}
\paragraph{Statement PL-02.}
Linearization preserves symmetric powers and their opposite-series exterior
operations, but the subset operations need not be special. Moreover, the
opposite-series structure itself need not be special: $G=C_2\times C_2$ gives
an explicit counterexample. Thus the free special $\lambda$-ring universal
property cannot simply be applied to these Burnside structures.

\begin{proof}
The degree-$i$ monomial basis of $\operatorname{Sym}^i\mathbb C[X]$ is indexed
by multisets of size $i$ in $X$, and the $G$-actions agree. Thus
$L(\sigma^i(X))=[\operatorname{Sym}^i\mathbb C[X]]$.
Disjoint unions give the addition law and extend this identity to virtual
classes. Applying $L$ to inverse series proves compatibility with
$\lambda^{\mathrm{op}}$. By contrast, in the example PL-01,
$L(\epsilon^2(x))=1$ whereas $\bigwedge^2\mathbb C[E]=\operatorname{sgn}$.

In every special $\lambda$-ring the second Adams operation
$\psi^2(z)=z^2-2\lambda^2(z)$ is multiplicative. At order two this follows
directly from the universal identity
\[
 \lambda^2(uv)=u^2\lambda^2(v)+v^2\lambda^2(u)
                       -2\lambda^2(u)\lambda^2(v).
\]
For the subset pre-structure on $B(C_2)$ one has
$\psi^2(x)=2x-2$, so
$\psi^2(x)^2=4$, while additivity gives
$\psi^2(x^2)=\psi^2(2x)=4x-4$. Their $C_2$-fixed-point marks are
$4$ and $-4$, so they differ.

For the opposite-series structure take $G=C_2\times C_2$, choose distinct
subgroups $H,K$ of order two, and let $a=[G/H]$, $b=[G/K]$.
Both sets have two elements and their actions factor through $C_2$. Hence
$a^2=2a$, $\sigma^2(a)=a+1$, and
$\psi^2(a)=2\sigma^2(a)-a^2=2$; likewise $\psi^2(b)=2$.
The product $(G/H)\times(G/K)$ is the regular four-element $G$-set $G$,
since its stabilizer is $H\cap K=1$. Write $r=[G]$. A $G$-fixed multiset on
the transitive set $G$ has constant multiplicity at all four points, so no
multiset of size two is fixed. Consequently
\[
 \Phi_G(\psi^2(r))=
  2\Phi_G(\sigma^2(r))-\Phi_G(r^2)=0,
 \qquad
 \Phi_G(\psi^2(a)\psi^2(b))=4.
\]
Multiplicativity fails, proving the claim.
\end{proof}

\subsection{PL-03: the characteristic-zero meaning of scan 28}
\paragraph{Statement PL-03.}
Let $V$ be a complex vector space, regarded universally as a representation of
$GL(V)$. Give $V^{\otimes n}$ its commuting place-permutation action of
$\mathfrak S_n$. Put $\rho_i=[\mathbf1_{\mathfrak S_i}]$, and multiply the
$\rho_i$ by induction. For $\alpha=(\alpha_1,\ldots,\alpha_n)$ with
$\sum i\alpha_i=n$, put $\mu(\alpha)=(1^{\alpha_1}2^{\alpha_2}\cdots)$.
Then the formula on cote 133, scan 28 has the following correct form in the
representation Grothendieck group:
\[
 [V^{\otimes n}]=
 \sum_{\alpha}m_{\mu(\alpha)}[V]\,
          \rho_1^{\alpha_1}\cdots\rho_n^{\alpha_n}.
 \tag{PL-03a}
\]
Here $m_\mu[V]$ means the monomial symmetric function expressed as an
integral polynomial in the exterior powers of $V$; it need not be effective.
In degree two the identity is
\[
 [V^{\otimes2}]=\psi^2[V]\,\rho_2+[\bigwedge^2V]\,\rho_1^2,
 \qquad \psi^2[V]=[\operatorname{Sym}^2V]-[\bigwedge^2V].
 \tag{PL-03b}
\]
The coefficient of $\rho_1^n$ is $e_n[V]$, and that of $\rho_n$ is
$p_n[V]=\psi^n[V]$. These coefficient assertions use the induction ring in
characteristic zero, where the $\rho_i$ freely generate a polynomial ring.

\begin{proof}
Restrict $GL(V)$ to a diagonal torus with eigencharacters $x_1,\ldots,x_r$.
For each exponent vector $(b_1,\ldots,b_r)$ with sum $n$, the tensors having
exactly $b_i$ copies of the $i$th basis vector span a torus-weight space of
weight $x_1^{b_1}\cdots x_r^{b_r}$. As an $\mathfrak S_n$-module this weight
space is the permutation module on ordered words of that content, namely
\[
 \operatorname{Ind}_{\mathfrak S_{b_1}\times\cdots\times\mathfrak S_{b_r}}
              ^{\mathfrak S_n}\mathbf1.
\]
Group exponent vectors by their unordered positive parts $\mu$.
The sum of their torus weights is $m_\mu(x)$, while their common permutation
module is the induction product $\prod_i\rho_i^{\alpha_i}$.
This proves (PL-03a) on the torus. Characters of polynomial $GL(V)$ modules
are symmetric polynomials, with the Schur characters linearly independent,
so torus restriction is injective on their Grothendieck group. Thus it proves
the asserted identity there. By universality it is the corresponding
polynomial identity in exterior powers in characteristic zero.
Finally $m_{(1^n)}=e_n$, $m_{(n)}=p_n$, and
$m_{(2)}=e_1^2-2e_2$; these give the coefficient statements and (PL-03b).
\end{proof}

For $\dim V\geq2$, $\psi^2[V]$ is not the class of a genuine polynomial
$GL(V)$ representation: its irreducible decomposition has coefficient $-1$
on $\bigwedge^2V$. Thus the terms of this particular monomial/induction
expansion need not be direct summands individually. In contrast the full
Schur--Weyl expansion
$V^{\otimes n}=\bigoplus_\lambda S_\lambda V\otimes[\lambda]$
is a direct-sum decomposition. Replacing one basis by another is exactly
where cancellation enters; the archive formula is not a positive combinatorial
formula for arbitrary plethysm multiplicities. This also gives the concrete
counterexample $\psi^2(s_1)=p_2=s_2-s_{11}$ to the inference that Adams
operations or Frobenius lifts must preserve Schur positivity.

\subsection{PL-04: the modular qualification cannot be omitted}
\paragraph{Statement PL-04.}
For the exact Grothendieck groups of finite-dimensional representations over
$k=\mathbb F_2$, the induction ring
$\bigoplus_nG_0(k[\mathfrak S_n])$ is not a polynomial ring freely generated
by the trivial classes $\rho_n$. Indeed $\rho_1^2=2\rho_2$.

\begin{proof}
The induction product $\rho_1^2$ is the regular $k[C_2]$-module.
Writing its generator as $g$, its invariant line is $k(1+g)$. In
characteristic two both this line and the one-dimensional quotient are
trivial modules. The short exact sequence therefore gives
$[k[C_2]]=2[k]$ in $G_0$, which is the asserted relation. A polynomial ring
on formal generators has no such relation, so the natural map from that
polynomial ring is not injective.
\end{proof}

The split Grothendieck group is a different object: the regular module here
is indecomposable, and the displayed short exact sequence does not split.
This distinction must be fixed before answering any question about arbitrary
local coefficient rings. The supplied study discusses scan 29, which was not
itself attached; PL-04 checks the quoted mathematical assertion, not that
page's wording. It does not classify all local rings for which a chosen
representation Grothendieck ring has polynomial generators.

\subsection{PL-05: positive sets cannot supply every positive representation}
\paragraph{Statement PL-05.}
An effective $G$-set linearizes to a character nonnegative on every element
of $G$. In particular, for $n\geq2$ the sign representation of $\mathfrak S_n$
is not the linearization of an effective finite $\mathfrak S_n$-set, even
though its class is positive in the representation semiring. Consequently a
construction that uses only effective sets, disjoint union, Cartesian product,
and linearization cannot produce all Schur-positive representations.

\begin{proof}
The permutation matrix of $g$ on $\mathbb C[X]$ has trace $|X^g|\geq0$.
The sign character has value $-1$ on a transposition, which rules out such
an $X$. Disjoint union and Cartesian product preserve the class of effective
sets, so adding these constructions does not remove the obstruction.
\end{proof}

This is a limitation on that proposed positive-set realization, not a
nonexistence theorem for positive formulas for Kronecker coefficients.
Restriction along the diagonal $G\to G\times G$ turns the external tensor
product into the internal one; cote 2, scan 17 writes exactly this operation.
For $G=\mathfrak S_n$ it is the operation whose irreducible multiplicities
are Kronecker coefficients. Nothing on the attached page identifies their
irreducible multiplicity spaces with explicitly countable sets.

\subsection{PL-06: the exact extra datum needed for the Chow obstruction}
Write $\widehat{\operatorname{Ch}}_n(W)\subset\operatorname{Sym}^nW$ for
the affine cone of products of $n$ linear vectors. The known
Hermite--Hadamard--Howe map is
\[
 h_{d,n}:\operatorname{Sym}^d(\operatorname{Sym}^nW^*)
       \longrightarrow\operatorname{Sym}^n(\operatorname{Sym}^dW^*).
\]
Its kernel is the degree-$d$ Chow ideal (Landsberg, attached survey,
Proposition 7.6). Its target is the degree-$d$ piece of the normalization's
coordinate ring (ibid., Propositions 7.11--7.12). If $\dim W=n$, the sum
of these targets is also the polynomial-representation part of the orbit
coordinate ring (ibid., \S7.5, using Theorem 3.7). Hence it is necessary
to distinguish orbit, orbit closure, and normalization.

\paragraph{Statement PL-06.}
Let $a_\lambda$ and $b_\lambda$ be the multiplicities of $S_\lambda W^*$
in the source and target of $h_{d,n}$. There is a linear map
$H_\lambda:\mathbb C^{a_\lambda}\to\mathbb C^{b_\lambda}$ such that
\[
 m_{\widehat{\operatorname{Ch}}_n,d}(\lambda)=\operatorname{rank}H_\lambda,
 \qquad
 m_{I_d}(\lambda)=a_\lambda-\operatorname{rank}H_\lambda.
 \tag{PL-06}
\]
In particular $0\leq m_{\widehat{\operatorname{Ch}}_n,d}(\lambda)
\leq\min(a_\lambda,b_\lambda)$. Knowledge of the two complete lists of
plethysm multiplicities does not determine these ranks for an arbitrary
equivariant map with that source and target. Even for the natural map,
the equality of source and target characters when $d=n$ does not force
injectivity.

\begin{proof}
Complete reducibility decomposes the two modules into
$\bigoplus_\lambda S_\lambda W^*\otimes\mathbb C^{a_\lambda}$ and
$\bigoplus_\lambda S_\lambda W^*\otimes\mathbb C^{b_\lambda}$.
Schur's lemma identifies an equivariant map with the sum of
$1_{S_\lambda W^*}\otimes H_\lambda$. Its image and kernel on each
summand have the multiplicities in (PL-06). Proposition 7.6 identifies its
image with the coordinate-ring quotient, proving the formula.
When $a_\lambda,b_\lambda>0$, the same source and target admit a zero map
and a rank-one multiplicity-space map, so their characters alone do not
force a rank. For the natural $h_{5,5}$ on $W=\mathbb C^5$, Landsberg's Proposition 7.9(1),
attributed there to Ikenmeyer--Mkrtchyan, supplies a nonzero kernel; for
example the module labelled $(14,7,2,2)$ occurs in it. Yet source and
target are the same module. This proves the final assertion by the cited
counterexample.
\end{proof}

Foulkes' multiplicity comparison and injectivity of a particular
Hermite--Hadamard--Howe map are therefore different assertions. A proof of
ordinary Schur positivity, or even an oracle giving all plethysm coefficients,
still needs control of $H_\lambda$ to produce these orbit-closure multiplicities.
For general GCT orbit closures the analogous missing datum is restriction
of functions from the ambient representation to the closure, including
boundary extension. This is a precise remaining gate, not a proof that a
future purity argument cannot control those maps.

\subsection{PL-07: the tensor-power bridge does not give saturation}
\paragraph{Statement PL-07.}
Kronecker coefficients fail the naive saturation implication
\[
 g(N\lambda,N\mu,N\nu)>0\ \Longrightarrow\ g(\lambda,\mu,\nu)>0.
\]
Specifically,
\[
 g((1,1),(1,1),(1,1))=0,\qquad
 g((2,2),(2,2),(2,2))=1.
\]

\begin{proof}
For $\mathfrak S_2$ the first invariant space is the invariant subspace of
the sign representation cubed, again sign, and is zero.
For $\mathfrak S_4$, let $M$ be its permutation representation on the
three partitions of four letters into two unordered pairs, and let
$U=M/\mathbf1$. On the conjugacy classes
$1,(12),(12)(34),(123),(1234)$, of sizes $1,6,3,8,6$, the numbers of
fixed pairings are $3,1,3,0,1$. Thus $U$ has character
$(2,0,2,-1,0)$. Its character has norm
$(4+3\cdot4+8)/24=1$, so $U$ is irreducible. The hook-length dimensions
of the five Specht modules for $\mathfrak S_4$ are $1,3,2,3,1$
(Macdonald, Chapter I, (7.6)(ii), and Section 5, Example 2), so it is
the module labelled $(2,2)$. The dimension of its tensor-cube invariant
space is the average of the cubed character:
\[
 \frac{1}{24}(1\cdot2^3+6\cdot0^3+3\cdot2^3+8\cdot(-1)^3+6\cdot0^3)
 =1.
\]
This proves both values.
\end{proof}

Ordinary representation positivity and the entire characteristic-zero
$\lambda$-ring dictionary already hold in this counterexample. Consequently
they cannot imply naive saturation. This leaves untouched carefully
qualified saturation-index hypotheses in GCT, whose exact quantifiers must
be considered separately.

\subsection{PL-08: two quadratic qualifications on scans 6 and 8}
For comparison, the usual Grothendieck--Witt ring of a scheme over a field
of characteristic different from two is known to be a special
$\lambda$-ring (Zibrowius, \emph{The $\gamma$-filtration on the Witt ring of
a scheme}, Theorem 2). The analogous statement for symmetric representations
of an affine algebraic group is Zibrowius, \emph{Symmetric representation
rings are $\lambda$-rings}, main theorem, p.~1056. These results concern
the Grothendieck--Witt definitions in those works; one must not silently
substitute the split group completion on an arbitrary locally ringed topos.
\paragraph{Statement PL-08.}
Exterior powers of underlying modules do not define rank-compatible
$\lambda$ operations on the group of strictly nondegenerate quadratic forms
in the generality proposed on cote 133, scan 6, when $2$ is not invertible.
Also, even when $2$ is invertible, the usual exterior-power operations on
Grothendieck--Witt groups need not descend to Witt groups, answering the
question on scan 8 negatively in that generality.

\begin{proof}
Over $\mathbb F_2$ the polar bilinear form of a quadratic form is alternating.
A nonsingular alternating form has even dimension: if $v\ne0$,
nonsingularity supplies $w$ with nonzero pairing, their span is a
nondegenerate plane, and induction on its orthogonal complement proves
the assertion. The hyperbolic quadratic plane $q(x,y)=xy$ has nonsingular
polar form. Its second exterior power as a vector space has dimension one,
so it cannot carry a strictly nondegenerate quadratic form. More generally
all ranks in the Grothendieck group generated by such forms are even,
whereas rank compatibility would require its second exterior operation
to have rank $\binom22=1$. This is a contradiction.

For the descent question work over $\mathbb R$, with symmetric bilinear
forms and the usual exterior-power form. The hyperbolic plane
$h=\langle1\rangle+\langle-1\rangle$ is zero in the Witt group.
Its exterior square is $\langle-1\rangle$. This is nonzero in the Witt
group: signature is additive, invariant under isometry, and zero on every
hyperbolic form, but its value on this line is $-1$. Thus a descended
operation would send the zero Witt class both to zero (from the zero
form) and to a nonzero class (from $h$). It is not well-defined.
\end{proof}

\subsection{PLE-F5: an invariant-ring realization of the obstructing ray}
\paragraph{Statement PLE-F5.}
Let $\zeta$ be a primitive sixth root of unity, and let $\Gamma$ act on
$\mathbb C[x,y]$ through the substitutions
\[
 a(x,y)=(\zeta x,\zeta^{-1}y),\qquad b(x,y)=(iy,ix).
\]
Then $\Gamma$ has order $12$, and
\[
 \mathbb C[x,y]^\Gamma\cong
 \mathbb C[U,V,Z]/(Z^2-UV^2-4U^4),
 \quad U=x^2y^2,\quad V=x^6-y^6,\quad Z=xy(x^6+y^6).
 \tag{PLE-F5a}
\]
Every invariant has even ordinary degree. Dividing degrees by two gives
$\deg U=2$, $\deg V=3$, $\deg Z=4$, and its Hilbert series and Hilbert
function are
\[
 \operatorname{Hilb}(t)=\frac{1+t^4}{(1-t^2)(1-t^3)},\qquad
 \dim\bigl(\mathbb C[x,y]^\Gamma\bigr)_s
      =\left\lfloor\frac{s}{3}\right\rfloor+\mathbf1_{2\mid s}.
 \tag{PLE-F5b}
\]
Consequently this function agrees with the plethysm/Kronecker ray
$\phi(s)$ in Kahle--Micha\l ek, \emph{Obstructions to combinatorial formulas
for plethysm}, Theorem 5 and the paragraph after Remark 13.
The dimension comparison here is strengthened in PLE-F8, which
constructs an explicit graded-algebra isomorphism with the
highest-weight multiplicity algebra of the cubic-plethysm ray.
No compatible Kronecker-algebra identification is asserted.
The invariant ring is normal, but its degree support
$\{0,2,3,4,\ldots\}$ is not saturated in the group $\mathbb Z$ it generates.

\begin{proof}
One checks $a^6=1$, $b^2=a^3$, and $bab^{-1}=a^{-1}$.
Every group element therefore has the form $a^j$ or $a^jb$, $0\leq j<6$.
The former substitutions are diagonal, the latter off-diagonal, and all
twelve are distinct, proving the order.
A monomial $x^py^q$ is invariant under $a$ exactly when $p-q$ is divisible
by six. If $p\geq q$, it is $(xy)^q(x^6)^{(p-q)/6}$; if $q\geq p$ use
$y^6$ instead. Hence
\[
 \mathbb C[x,y]^{\langle a\rangle}
   =\mathbb C[A,B,C]/(AB-C^6),
 \quad A=x^6,\ B=y^6,\ C=xy.
\]
To justify that there is no further relation, reduce a polynomial modulo
$AB-C^6$ to a linear combination of monomials $A^rC^q$ ($r,q\geq0$)
and $B^rC^q$ ($r\geq1,q\geq0$); their images are distinct monomials in
$x,y$, so they are linearly independent.
The remaining involution sends $(A,B,C)$ to $(-B,-A,-C)$.
Putting $V=A-B$, $D=A+B$ gives
\[
 \mathbb C[A,B,C]/(AB-C^6)
   =\mathbb C[C,V,D]/(D^2-V^2-4C^6),
\]
where the involution fixes $V$ and negates $C,D$.
Because the relation is monic in $D$, every class uniquely has the form
$f(C,V)+Dg(C,V)$. It is invariant exactly when $f$ is even in $C$ and
$g$ is odd in $C$. The invariant ring therefore is the free
$\mathbb C[U=C^2,V]$-module with basis $1,Z=CD$; its multiplication is
specified by
$Z^2=C^2D^2=UV^2+4U^4$. This proves (PLE-F5a).

The two basis elements have half-degrees zero and four, proving the
Hilbert series in (PLE-F5b). Independently,
\[
 \sum_{s\geq0}\left\lfloor\frac{s}{3}\right\rfloor t^s
       =\frac{t^3}{(1-t)(1-t^3)},\qquad
 \sum_{s\geq0}\mathbf1_{2\mid s}t^s=\frac1{1-t^2}.
\]
Adding and cancelling $1-t$ gives the same Hilbert series. Finally the
residue classes $s=0,1,\ldots,5\pmod6$ give the periodic correction
$(3,-1,1,0,2,-2)/3$ in $\phi(s)=(s+r(s))/3$, which is the cited ray.

Normality has a direct proof. Put $A=\mathbb C[x,y]$ and $B=A^\Gamma$.
If $f\in\operatorname{Frac}(B)$ is integral over $B$, the same monic
equation shows that it is integral over $A$. Since $A$ is integrally
closed in $\operatorname{Frac}(A)$, we have $f\in A$. As a fraction of
elements of $B$, it is fixed by $\Gamma$, so $f\in B$. Thus $B$ is
normal. Its displayed Hilbert function is zero in degree one and positive
in every degree at least two. Its support generates $\mathbb Z$ because
it contains two and three, but is not saturated because $2\cdot1$ belongs
to the support and $1$ does not.
\end{proof}

The invariant-ring realization shows directly why a positive Hilbert
function and an ordinary homogeneous Ehrhart function are distinct
requirements. The ring is normal and has a positive finite free-module
basis over $\mathbb C[U,V]$, but the associated ray still violates
homogeneous Ehrhart reciprocity and its degree support is not saturated.
These are properties of this weighted grading; normality of a toric
semigroup algebra with its full semigroup grading is a different assertion.
No RH or GRH assumption was used in this calculation.
Positive combinatorial formulas already cover all plethysm coefficients
with two-row output and the corresponding equal-rectangle Kronecker family:
Pak--Panova--Swanson, \emph{A combinatorial interpretation for certain
plethysm and Kronecker coefficients}, arXiv:2511.02312v3 (July 2026),
Theorems 1.1--1.2 and Lemma 2.1. Thus PLE-F5 is an explicit invariant-ring
realization of one ray, not a claim to a first positive formula for that ray.

\subsection{PLE-F6: exactly which rescalings repair homogeneous counting}
\paragraph{Statement PLE-F6.}
For an integer $q\ge1$, define
$\phi_q(s)=\phi(qs)$ for integers $s\ge0$, with $\phi$ as in PLE-F5.
There is a rational polytope $Q$ such that
$\phi_q(s)=\#(sQ\cap\mathbb Z^D)$ for every $s\ge0$ if and only if
$\gcd(q,6)>1$. More precisely, when $q$ is even the rational interval
$Q=[0,q/3]$ works. When $q$ is odd and divisible by three, the interval
$Q=[1/2,q/3+1/2]$ works. Hence two is the smallest positive rescaling
that repairs homogeneous Ehrhart counting on this ray.

\begin{proof}
If $q$ is even, the formula in PLE-F5 gives
\[
 \phi_q(s)=\left\lfloor\frac{qs}{3}\right\rfloor+1
 =\#\bigl([0,qs/3]\cap\mathbb Z\bigr),
\]
including $s=0$. If $q$ is odd and divisible by three, put $k=q/3$.
Then
\[
 \phi_q(s)=ks+\mathbf1_{2\mid s}
 =\left\lfloor(k+1/2)s\right\rfloor
  -\left\lceil s/2\right\rceil+1,
\]
which is the claimed interval count, again including $s=0$.

It remains to exclude odd $q$ not divisible by three. Suppose a rational
polytope $Q$ gave the count. Its dimension would be one. Indeed,
dimension zero gives at most one point. In dimension $d\ge2$, choose
$d+1$ affinely independent rational vertices and clear their denominators
by $M$, obtaining integral vertices $v_0,\ldots,v_d$ in $MQ$. The distinct
integral points
$tv_0+\sum_{j=1}^d k_j(v_j-v_0)$ with $k_j\ge0$ and
$\sum_jk_j\le t$ all belong to $MtQ$. There are
$\binom{t+d}{d}$ of them, contradicting
$\phi_q(Mt)=qMt/3+O(1)$.
The quasi-polynomial extension of the formula for $\phi_q$ to negative
integers therefore obeys one-dimensional Ehrhart reciprocity, as stated
in Kahle--Micha\l ek, Remark 6. At $s=1$ its interior count would be
\[
 \#\bigl(\operatorname{relint}(Q)\cap\mathbb Z^D\bigr)
 =-\phi_q(-1)
 =-\left\lfloor-\frac q3\right\rfloor
 =\left\lceil\frac q3\right\rceil.
\]
But its closed count is
$\phi_q(1)=\lfloor q/3\rfloor$ since $q$ is odd. These two integers
differ by one since $3\nmid q$, contradicting inclusion of the relative
interior in $Q$. No rational homogeneous model exists in this case.
The two constructive cases and this exclusion exhaust all $q\ge1$.
\end{proof}

Kahle--Micha\l ek, Question 15 and its following paragraph, already
consider rescaling a stretching function into an Ehrhart function and
note a positive result in degree one. PLE-F6 records the exact allowed
rescalings for this particular ray. It does not restore the original
$q=1$ instance of PH1, and it does not provide a uniform polynomial bound
for rescaling arbitrary plethysm data.

\subsection{PLE-F7: pure Tate cohomology on the same obstructing cone}
\paragraph{Statement PLE-F7.}
Let $q$ be a prime power with $q\equiv1\pmod{12}$ and let
\[
 X=\operatorname{Spec}
 \mathbb F_q[U,V,Z]/(Z^2-UV^2-4U^4),
 \qquad \deg(U,V,Z)=(2,3,4).
\]
This is a normal affine surface with the Hilbert function $\phi$ of PLE-F5.
For every prime $\ell\nmid6q$, its geometric compactly supported
cohomology is
\[
 H_c^j(X_{\overline{\mathbb F}_q},\mathbb Q_\ell)
 \cong
 \begin{cases}\mathbb Q_\ell(-2)&j=4,\\0&j\ne4.\end{cases}
 \tag{PLE-F7}
\]
The identification respects Frobenius, so this cohomology is even,
Tate, and pure of the expected weight. Nevertheless the Hilbert function
is not a homogeneous rational Ehrhart function and the degree support
is not saturated. With the perverse normalization
$IC_X=j_{!*}\mathbb Q_{\ell,X_{\rm sm}}[2]$, one also has
$IC_X\simeq\mathbb Q_{\ell,X}[2]$. Consequently
$H_c^j(X_{\overline{\mathbb F}_q},IC_X)$ is zero except for
$j=2$, where it is $\mathbb Q_\ell(-2)$.
In fact the displayed hypersurface has
$q^{2r}$ points over $\mathbb F_{q^r}$ for every $r\ge1$ for every odd
$q$, even without $q\equiv1\pmod{12}$.

\begin{proof}
For $q\equiv1\pmod{12}$, choose $\zeta,i\in\mathbb F_q$ of orders
six and four. The substitutions in PLE-F5 define the same group $\Gamma$
of order twelve. Every step in the invariant-ring calculation in that
proof is valid over $\mathbb F_q$: the cyclic weights are distinct and
two is invertible. Thus $X=\mathbb A^2/\Gamma$ with the displayed
coordinate ring, free over $\mathbb F_q[U,V]$ with basis $1,Z$.
The same integral-closure argument proves normality, and the free basis
proves the same Hilbert function. In particular $X$ has dimension two.

We give the cohomological comparison explicitly, so no identification
between graded multiplicities and cohomological degrees is implicit.
The quotient map $\pi:\mathbb A^2\to X$ is finite: each element $a$ of
the polynomial ring satisfies the monic equation
$\prod_{g\in\Gamma}(T-ga)$ with invariant coefficients, and the ring
is generated by its two integral coordinate functions.
Each geometric fiber is a single $\Gamma$-orbit as a set. To see the
last assertion, two distinct finite orbits in $\overline{\mathbb F}_q^2$
can be separated by a polynomial that is zero on the first and one on
the second, by the Chinese remainder theorem for their distinct maximal
ideals. Averaging that polynomial over $\Gamma$ produces an invariant
with the same values, contradicting equality of their quotient points.
Surjectivity follows from the integral extension. The same explicit
invariant presentation is valid after extending the field.

Put $\Lambda_n=\mathbb Z/\ell^n\mathbb Z$. For a finite morphism,
the stalk of $\pi_*\Lambda_n$ is the space of functions on the finite
geometric fiber and $R^j\pi_*\Lambda_n=0$ for $j>0$
(Stacks Project, Proposition 59.55.2, tag \texttt{03QP}). Since the fibers
are single orbits, the unit map identifies
\[
 \Lambda_n\simeq(\pi_*\Lambda_n)^\Gamma
 =\operatorname{im}\left(e:\pi_*\Lambda_n\to\pi_*\Lambda_n\right),
 \qquad e=\frac1{12}\sum_{g\in\Gamma}g.
\]
Here $12$ is invertible in $\Lambda_n$, and $e^2=e$, so this is a direct
summand of sheaves. Finite maps are proper, and the composition theorem
for cohomology with proper support (Milne, \emph{Lectures on \'Etale
Cohomology}, version 2.21, Proposition 18.4(d)) therefore gives
\[
 H_c^j(X_{\overline{\mathbb F}_q},\Lambda_n)
 =eH_c^j(X_{\overline{\mathbb F}_q},\pi_*\Lambda_n)
 =H_c^j(\mathbb A^2_{\overline{\mathbb F}_q},\Lambda_n)^\Gamma.
\]
Affine-space acyclicity and Poincar\'e duality give only
$H_c^4(\mathbb A^2,\Lambda_n)=\Lambda_n(-2)$
(Milne, Example 16.3 and Theorem 24.1).
The trace-normalized class of the origin generates the top group after
the Tate twist. Every element of $\Gamma$ fixes the origin, hence fixes
this generator. Taking the inverse limit in $n$ and tensoring with
$\mathbb Q_\ell$ proves (PLE-F7). All maps are defined over
$\mathbb F_q$, so the assertion is Frobenius-equivariant; geometric
Frobenius acts by $q^2$ on $\mathbb Q_\ell(-2)$.

We verify the intersection-complex assertion separately. The partial
derivatives of $Z^2-UV^2-4U^4$ are
$2Z,-2UV,-V^2-16U^3$. Since the characteristic is not two, their
common zero on $X$ is only the origin $o$. Thus
$j:X\setminus\{o\}\hookrightarrow X$ is the smooth locus.
The fiber of $\pi$ over $o$ consists only of the origin in $\mathbb A^2$:
$U=x^2y^2=0$ and $V=x^6-y^6=0$ force $x=y=0$.
Apply the preceding finite-pushforward and averaging argument to the
localization triangles for the pairs $(X,\{o\})$ and
$(\mathbb A^2,\{0\})$. It gives
\[
 R\Gamma_{\{o\}}(X_{\overline{\mathbb F}_q},\mathbb Q_\ell)
 \simeq
 R\Gamma_{\{0\}}(\mathbb A^2_{\overline{\mathbb F}_q},
                 \mathbb Q_\ell)^\Gamma
 \simeq \mathbb Q_\ell(-2)[-4].
\]
The second isomorphism is purity for the smooth pair
$(\{0\},\mathbb A^2)$ (Milne, Theorem 16.1); the generator is again
the class of the origin and is fixed by $\Gamma$.
Let $i:\{o\}\hookrightarrow X$ and $K=\mathbb Q_{\ell,X}[2]$.
Since a geometric point has no higher cohomology, the local calculation
and the constant stalk give
\[
 i^*K=\mathbb Q_\ell[2],\qquad
 i^!K=\mathbb Q_\ell(-2)[-2].
\]
On the smooth open locus $K$ is a local system shifted by its dimension;
at $o$ its stalk and costalk satisfy the middle-perversity conditions.
Thus $K$ is perverse (BBD, \emph{Faisceaux pervers},
(4.0.1)--(4.0.2)). Moreover $H^0(i^*K)=H^0(i^!K)=0$.
The intermediate-extension criterion (BBD, Corollaries 1.4.24--1.4.25)
therefore identifies $K$ with $j_{!*}\mathbb Q_{\ell,X\setminus\{o\}}[2]$:
it has no nonzero perverse subobject or quotient supported at $o$.
The identification restricts to the identity on the smooth locus;
uniqueness of intermediate extension makes it Frobenius-equivariant.
The stated intersection-cohomology groups now follow by this explicit
shift, independently of the point count below.

There is also an elementary, independent point count. Let $Q$ be any odd
prime power and let $\chi$ be the quadratic character of $\mathbb F_Q$,
extended by $\chi(0)=0$. For $c\ne0$,
\[
 \sum_{v\in\mathbb F_Q}\chi(v^2+c)=-1.
\]
Indeed the number of pairs $(v,w)$ with $w^2-v^2=c$ is $Q-1$,
because $(w-v)(w+v)=c$ is parametrized by an arbitrary nonzero first
factor and two is invertible; counting first by $v$ gives $Q$ plus the
displayed sum. Therefore
\[
 \#X(\mathbb F_Q)
 =\sum_{u,v\in\mathbb F_Q}\bigl(1+\chi(uv^2+4u^4)\bigr)
 =Q^2-\sum_{u\ne0}\chi(u)=Q^2.
\]
The $u=0$ character contribution is zero; for $u\ne0$ use
$uv^2+4u^4=u(v^2+4u^3)$ and the preceding identity. Taking $Q=q^r$
proves the stated point counts. The Hilbert-function and support
failures are those already proved in PLE-F5 and PLE-F6.
\end{proof}

Thus the pure Tate compactly supported cohomology of this affine
quotient, including its intersection cohomology, even with normality
and a positive finite free-module Hilbert description, cannot supply
the missing homogeneous counting assertion.
The ordinary compactly supported cohomology in PLE-F7 has only one
nonzero degree, and the intersection complex is the constant sheaf
shifted by two. The grading whose pieces have dimensions $\phi(s)$ is
different data. PLE-F8 identifies the graded algebra with actual
plethysm multiplicities and supplies a positive multiplicative basis.
It does not identify those graded pieces with cohomological groups
carrying all the additional structure in the categorical criterion
P-F6. This example does not refute a proposed construction with
those additional properties.

\subsection{Archive locator inventory and the scope of the bridge}
\begin{description}
\item[Cote 133, scan 6.] Quadratic Grothendieck--Witt groups, and explicit
questions about the universal product and iteration formulas. This is not
a proof that the quadratic group is a special $\lambda$-ring. The page's
warning about noninvertible $2$ matters; PL-08 gives a counterexample to
the unrestricted exterior-power assertion.
\item[Cote 133, scan 8.] The hyperbolic ideal and the question whether the
pre-$\lambda$ operations descend to its quotient. A question on the sheet
cannot be used as an affirmative theorem; PL-08 answers it negatively
for the standard operations over $\mathbb R$.
\item[Cote 133, scan 9.] Stiefel--Whitney operations and universal tensor
formulas; the assertion that the total class factors through the Witt
quotient is visibly deleted in the transcription.
\item[Cote 133, scan 19.] Subsets, multisets, their proposed comparison, and
an open product-formula question. PL-01--PL-02 give the needed equivariant
correction.
\item[Cote 133, scan 28.] Tensor powers with symmetric-group action, the
induction product, and monomial symmetric functions in the exterior powers.
PL-03 gives its characteristic-zero meaning; PL-04 explains why integral
or modular variants cannot silently use the same freeness.
\item[Cote 2, scan 17.] External tensor products, pullback, projection
formula, and restriction along the diagonal. This supplies the categorical
operation for the inner tensor product, not its numerical multiplicities.
\item[Cote 43, scan 37.] The ring of augmentation-one formal series,
convolution, $\lambda$ operations, and Frobenius--Verschiebung formulas in
characteristic $p$. The last map to $p$-typical Witt vectors carries a
question mark. Formal Frobenius structure is not coefficient positivity;
$p_2=s_2-s_{11}$ is an explicit distinction.
\item[Cote 119, scan 217 (site foliation 216).] A retrospective attribution
of $K$-theory and $\lambda$ operations to the Riemann--Roch work and their
links with symmetric-group operations; no plethysm positivity theorem is
stated on this sheet.
\item[Cote 12, scan 41.] A proposed determination of $\lambda$ operations
through finite operator groups, not a worked-out coefficient formula.
\item[Cote 16, scan 38.] Lefschetz operators, primitive parts, and
hyperplane restriction and pushforward identities. Its capital $\Lambda$ is a
Lefschetz operator, not a plethystic exterior-power operation.
\item[Cote 4, scan 7.] Crystalline/\'etale vanishing cohomology, torsion,
and unit Frobenius roots, with a correction marked false in the margin.
This excerpt does not identify any plethysm multiplicity with that
cohomology.
\item[Cote 79, scan 51.] Scalar torsor actions on triangulated finite
polyhedral structures. Its $\lambda$ is a scalar/action parameter.
\item[Cote 89, scan 18.] Reconstruction of a vector space from punctured
lines and wedge/torsor data. Its $\Lambda$ is a coefficient ring. Neither
of these last two sheets supplies $\lambda$-ring positivity.
\end{description}

Thus the attached pages contain substantial algebra of operations, including
the correct source of tensor, Schur, and plethystic constructions. They do
not state a positive lattice-point rule, a saturation theorem, or an
asymptotic determinant--permanent obstruction. The meaningful next target
is a realization of the relevant multiplicity spaces together with their
restriction maps and quantitative complexity bounds; the formal algebra
alone has the counterexamples above.
